\section{Описание}

\subsection{Общее устройство}
Поисковый робот --- программа \texttt{robot.py} на Python 3 (библиотеки
\texttt{requests}, \texttt{pymongo}, \texttt{PyYAML}). Вся логика --- один класс
\texttt{MusicCrawler}. Робот обкачивает те же два источника, что использовались
при сборке корпуса в первой лабораторной работе: Lyrics.ovh API и MusicBrainz
API. Хранилище --- MongoDB.

Робот не бродит по ссылкам: список URL для обкачки строится из метаданных
корпуса из первой лабораторной работы (файлы \texttt{*.json} в каталогах
\texttt{metadata}). Для каждого трека из метаданных собирается URL вида
\texttt{https://api.lyrics.ovh/v1/\{артист\}/\{название\}}, для каждой записи
MusicBrainz --- \texttt{https://musicbrainz.org/ws/2/recording/\{mbid\}?fmt=json}.
Таким образом, корпус из первой работы выступает seed-списком, а робот
переобкачивает те же документы и следит за их изменениями.

\subsection{Конфигурация}
Единственный аргумент робота --- путь до YAML-конфига. В нём три обязательные
секции: \texttt{db} (подключение к MongoDB), \texttt{logic} (задержка между
обкачкой, лимит документов на источник, интервал переобкачки, таймауты, число
повторных попыток) и \texttt{sources} (источники данных и пути к метаданным
корпуса). Дополнительно --- секции \texttt{robots\_txt} и \texttt{logging}.

\begin{lstlisting}[language=Python]
db:
  host: localhost
  port: 27017
  database: robot
  collection: documents

logic:
  delay_between_requests: 2.0
  max_documents_per_source: 15000
  refresh_interval_days: 7
  retry_attempts: 3
  timeout_seconds: 30

sources:
  lyrics_ovh:
    metadata_dir: "../lab1/lyrics_corpus/metadata"
    source_name: "Lyrics.ovh"
  musicbrainz:
    metadata_dir: "../lab1/musicbrainz_corpus/metadata"
    source_name: "MusicBrainz"
    respect_robots_txt: false   # /ws/2 is the official public API
\end{lstlisting}

\subsection{Схема документа в MongoDB}
Каждый обкачанный документ --- запись в коллекции \texttt{documents}:

\begin{tabular}{|l|p{10cm}|}
\hline
Поле & Описание \\ \hline
\texttt{url} & нормализованный URL (scheme + host + path, для MusicBrainz ---
без query-параметров); уникальный индекс \\ \hline
\texttt{html\_content} & \enquote{сырой} текст ответа (для API-источников ---
тело JSON-ответа как есть) \\ \hline
\texttt{source} & название источника (Lyrics.ovh / MusicBrainz) \\ \hline
\texttt{crawl\_timestamp} & дата обкачки, Unix time stamp \\ \hline
\texttt{content\_hash} & MD5-хэш содержимого, для проверки изменений \\ \hline
\texttt{title}, \texttt{content\_length}, \ldots & заголовок документа и
служебные поля (код ответа, Content-Type, Last-Modified) \\ \hline
\texttt{previous\_versions} & история предыдущих версий изменившихся
документов (хэш, дата, заголовок) \\ \hline
\end{tabular}

\subsection{Возобновление работы}
После каждого успешно сохранённого документа робот записывает контрольную
точку --- последний обработанный URL --- в отдельную коллекцию
\texttt{crawler\_checkpoints}. При запуске точка загружается, список URL
строится заново (детерминированно, из тех же метаданных), и робот продолжает
со следующего после контрольной точки документа. Остановка --- по
\texttt{Ctrl+C} (SIGINT/SIGTERM): робот завершает текущий документ, сохраняет
статистику и выходит.

\subsection{Переобкачка изменившихся документов}
Документ переобкачивается, если с прошлой обкачки прошло больше
\texttt{refresh\_interval\_days} дней. Скачанный контент сравнивается по
MD5-хэшу с сохранённым: если не изменился --- обновляется только
\texttt{crawl\_timestamp}; если изменился --- документ обновляется, а старая
версия попадает в \texttt{previous\_versions}.

\subsection{Особенности реализации}
При написании кода пришлось учесть особенности источников:

\begin{itemize}
    \item оба источника --- JSON API, а не HTML-сайты: робот принимает ответы
    с \texttt{Content-Type} \texttt{text/html} и \texttt{application/json},
    \enquote{сырой} текст ответа сохраняется как есть;
    \item robots.txt MusicBrainz содержит \texttt{Disallow: /ws}, под который
    попадает официальное публичное API \texttt{/ws/2}; поэтому проверка
    robots.txt отключается отдельно для этого источника (параметр
    \texttt{respect\_robots\_txt}), для остальных источников она включена;
    \item в метаданных корпуса есть дубли (одинаковые пары артист---название),
    поэтому список URL дедуплицируется перед обкачкой;
    \item сетевые сбои: повторные попытки запроса (3 раза, с backoff), таймаут
    30~с, пауза 2~с между запросами, чтобы не долбить API.
\end{itemize}

\subsection{Статистика запуска}
Статистика приведена для целевого размера корпуса --- 30000 документов (по
15000 из каждого источника); значения экстраполированы по контрольному прогону
(108 документов за 292~с, прерывание по \texttt{Ctrl+C}).

\begin{tabular}{|l|c|}
\hline
Показатель & Значение \\ \hline
Количество документов & 30000 \\ \hline
Скорость обкачки & $\sim$0{,}37 док/с (задержка 2~с + время запроса) \\ \hline
Время полной обкачки & $\sim$22--23 часа \\ \hline
Средний размер сырого ответа & $\sim$1 КБ (тексты песен $\sim$1{,}5 КБ,
записи $\sim$0{,}3 КБ) \\ \hline
Объём сырых данных в БД & $\sim$30 МБ \\ \hline
Ошибки обкачки & $\sim$0\% (в контрольном прогоне --- 0 из 108) \\ \hline
Интервал переобкачки & 7 дней \\ \hline
\end{tabular}

Возобновление работы проверено: после остановки и повторного запуска робот
загрузил контрольную точку и продолжил со следующего документа
(\texttt{Продолжаем с позиции 2} в логе).

\pagebreak

\section{Исходный код}
Программа \texttt{robot.py} состоит из класса \texttt{MusicCrawler}. Основные
блоки: построение списка URL из метаданных, обкачка одного URL
(\texttt{crawl\_url}), сохранение документа с проверкой изменений
(\texttt{process\_document}), контрольные точки
(\texttt{save\_checkpoint}/\texttt{load\_checkpoint}) и основной цикл
(\texttt{run}).

Обкачка одного URL: нормализация, проверка robots.txt, запрос с таймаутом,
проверка Content-Type, вычисление хэша:

\begin{lstlisting}[language=Python]
def crawl_url(self, url, source_name, respect_robots_txt=True):
    normalized_url = self.normalize_url(url, source_name)
    if respect_robots_txt and not self.check_robots_txt(url, source_name):
        return None
    response = self.session.get(url, timeout=...)
    response.raise_for_status()
    content = response.text
    content_type = response.headers.get('Content-Type', '')
    if 'text/html' not in content_type and 'json' not in content_type:
        return None
    return {
        'url': normalized_url,
        'html_content': content,
        'content_hash': self.calculate_content_hash(content),
        'source': source_name,
        'crawl_timestamp': int(time.time()),
        ...
    }
\end{lstlisting}

Сохранение документа: сравнение хэша с сохранённым, обновление только
изменившихся, история предыдущих версий:

\begin{lstlisting}[language=Python]
def process_document(self, document):
    existing = self.collection.find_one({'url': document['url']})
    if existing:
        if existing.get('content_hash') == document['content_hash']:
            # unchanged: only refresh the crawl date
            self.collection.update_one(
                {'_id': existing['_id']},
                {'$set': {'crawl_timestamp': document['crawl_timestamp']}})
        else:
            # changed: old version goes to history, document updated
            document['previous_versions'] = existing.get(
                'previous_versions', []) + [{
                    'content_hash': existing.get('content_hash'),
                    'crawl_timestamp': existing.get('crawl_timestamp')}]
            self.collection.update_one(
                {'_id': existing['_id']}, {'$set': document})
    else:
        self.collection.insert_one(document)
\end{lstlisting}

Контрольная точка --- последний обработанный URL (коллекция
\texttt{crawler\_checkpoints}, одна запись на робота):

\begin{lstlisting}[language=Python]
def save_checkpoint(self, url):
    self.db['crawler_checkpoints'].update_one(
        {'crawler_id': 'music_crawler'},
        {'$set': {'last_url': url, 'timestamp': time.time()}},
        upsert=True)
\end{lstlisting}

Основной цикл: пропуск ещё свежих документов, обкачка, задержка между
запросами:

\begin{lstlisting}[language=Python]
for i in range(start_index, len(all_urls)):
    if not self.running:          # stop on Ctrl+C
        break
    url_info = all_urls[i]
    existing = self.collection.find_one({'url': url_info['url']})
    if existing and not self.should_recrawl(existing):
        continue                  # too early to re-crawl
    document = self.crawl_url(url_info['url'], ...)
    if document:
        self.process_document(document)
        self.save_checkpoint(url_info['url'])
    time.sleep(self.config['logic']['delay_between_requests'])
\end{lstlisting}

\pagebreak
